Across 15 Fashion-MNIST training runs, the custom Transformer achieves the highest
mean test accuracy and macro-F1, 91.267\% and 91.234\%, narrowly ahead of the
custom CNN at 90.690\% and 90.726\%.

The comparison uses five architecture families, one persisted stratified split and
three seeds per configuration. Selected checkpoints are evaluated on the official
10,000-image test source, with mean and sample standard deviation reported across seeds.

The CNN offers close predictive performance with about 72.9\% fewer parameters
than the Transformer and lower observed training-loop cost. Separately measured
Windows CPU inference favors different trade-offs; the model choice depends on
quality, size and the intended runtime.

All classifiers use shared preprocessing and metrics, while their representations
range from flattened pixels to spatial feature maps, image rows and patch tokens.
The study reports validation/test metrics, exact parameter counts, learning curves,
T4 train-and-validation loop times, explicitly labeled CPU timing, confusion
matrices and selected successes/failures. Shirt and related upper-body categories
remain difficult across models. Small seed count, different training policies and
separate timing environments prevent a causal or universal ranking of model families.
Effective configurations, run/checkpoint identities, verification records and
reconstruction instructions accompany the editable report source.
